% 6th Grade Review Explanations — corrected slice: G1, H1–H7, I1–I3, J1–J5
% (source pages 23–33). Fragment: \input into a wrapper that loads mathdocs.sty.

\lesson{G. Distributive Property}

\question{G1}

\blockhead{Question}
Expand: $3(y + 6) = $

\checked{$3y + 18$}

\blockhead{What the question is asking}
This question asks you to rewrite $3(y + 6)$ without curved marks by multiplying
the outside number by each inside part.

\blockhead{Important words}
\begin{words}
  \item The \term{distributive property}says $a(b + c) = ab + ac$. You multiply
        the outside number by each part inside.
  \item To \term{expand}means write the expression without the parentheses by
        distributing.
\end{words}

\blockhead{Work the problem one tiny step at a time}
\begin{steps}
  \item The outside number is 3. The inside parts are $y$ and 6.
  \item Multiply 3 by the first inside part: $3 \times y = 3y$.
  \item Multiply 3 by the second inside part: $3 \times 6 = 18$.
  \item Add those products: $3y + 18$.
\end{steps}

\finalans{$3y + 18$}

\blockhead{Why this makes sense}
Check with a number: if $y = 2$, left side $3(2 + 6) = 3(8) = 24$. Right side
$3(2) + 18 = 6 + 18 = 24$. Same value.

\lesson{H. Equations, Inequalities, And Graphing}

\question{H1}

\blockhead{Question}
$x + 7 = 18$\qquad $x = $

\checked{$x = 11$}

\blockhead{What the question is asking}
This question asks you to find the number $x$ that makes both sides equal.

\blockhead{Important words}
\begin{words}
  \item An \term{equation}says two sides are equal.
  \item To \term{solve}means find the value of the variable.
  \item Use the \term{inverse operation}: undo $+7$ by subtracting 7 from both sides.
  \item To \term{substitute}means put the number back in for the letter and check
        both sides.
\end{words}

\blockhead{Work the problem one tiny step at a time}
\begin{steps}
  \item Goal: get $x$ alone on one side so both sides stay equal.
  \item Start with $x + 7 = 18$.
  \item The left side shows $x$ plus 7. The inverse of add 7 is subtract 7.
  \item Subtract 7 from both sides. Show the operation under the equals:
        \begin{center}
          $x + 7 = 18$\par\smallskip
          \underline{\hspace{0.6em}\textcolor{blue}{$-7$}\hspace{1.6em}\textcolor{blue}{$-7$}\hspace{0.6em}}
        \end{center}
  \item Left side: $x + 7 - 7$ leaves $x$.
  \item Right side leftover arithmetic: $18 - 7 = 11$.
  \item So $x = 11$.
  \item Check by substituting: put 11 back in for $x$. Compute $11 + 7 = 18$.
        Both sides match.
\end{steps}

\finalans{$x = 11$}

\blockhead{Why this makes sense}
Check: put 11 back in. $11 + 7 = 18$. Both sides match, so $x = 11$ is correct.

\question{H2}

\blockhead{Question}
$\dfrac{x}{4} = 6$\qquad $x = $

\checked{$x = 24$}

\blockhead{What the question is asking}
This question asks you to find the number $x$ that makes $\dfrac{x}{4} = 6$ true.

\blockhead{Important words}
\begin{words}
  \item $\dfrac{x}{4}$ means $x$ divided by 4.
  \item The inverse of divide by 4 is multiply by 4. Do that to both sides.
\end{words}

\blockhead{Work the problem one tiny step at a time}
\begin{steps}
  \item Goal: get $x$ alone on one side so both sides stay equal.
  \item Start with $\dfrac{x}{4} = 6$.
  \item The left side shows $x$ divided by 4. The inverse of divide by 4 is
        multiply by 4.
  \item Multiply both sides by 4. Show the operation under the equals:
        \begin{center}
          $\dfrac{x}{4} = 6$\par\smallskip
          \underline{\hspace{0.6em}\textcolor{blue}{$\times 4$}\hspace{1.6em}\textcolor{blue}{$\times 4$}\hspace{0.6em}}
        \end{center}
  \item Left side: $\dfrac{x}{4} \times 4$ leaves $x$.
  \item Right side leftover arithmetic: $6 \times 4 = 24$.
  \item So $x = 24$.
  \item Check by substituting: put 24 back in for $x$. Compute
        $\dfrac{24}{4} = 6$. Both sides match.
\end{steps}

\finalans{$x = 24$}

\blockhead{Why this makes sense}
Check: $\dfrac{24}{4} = 6$. Yes, both sides match.

\question{H3}

\blockhead{Question}
Solve: $x + 3 > 8$

\checked{$x > 5$}

\blockhead{What the question is asking}
This question asks you to find which numbers for $x$ make the left side of
$x + 3 > 8$ greater than the right.

\blockhead{Important words}
\begin{words}
  \item An \term{inequality}compares two sides with a symbol like $>$ (greater
        than) or $<$ (less than).
  \item $>$ means ``greater than'' (left side is bigger).
  \item To \term{isolate}$x$ means get the letter $x$ alone on one side.
  \item You can add or subtract the same number on both sides of an inequality,
        just like an equation.
\end{words}

\blockhead{Work the problem one tiny step at a time}
\begin{steps}
  \item Start with $x + 3 > 8$.
  \item Subtract 3 from both sides to isolate $x$. Show under the inequality:
        \begin{center}
          $x + 3 > 8$\par\smallskip
          \underline{\hspace{0.6em}\textcolor{blue}{$-3$}\hspace{1.6em}\textcolor{blue}{$-3$}\hspace{0.6em}}
        \end{center}
  \item Left side: $x + 3 - 3$ leaves $x$.
  \item Right side leftover arithmetic: $8 - 3 = 5$.
  \item So $x > 5$.
\end{steps}

\finalans{$x > 5$}

\blockhead{Why this makes sense}
Any number greater than 5 works. Example: if $x = 6$, then $6 + 3 = 9$ and
$9 > 8$ is true. If $x = 5$, then $5 + 3 = 8$, and $8 > 8$ is false, so 5 is not
included.

\question{H4}

\blockhead{Question}
Solve: $x - 4 < 2$

\checked{$x < 6$}

\blockhead{What the question is asking}
This question asks you to find which numbers for $x$ make the left side of
$x - 4 < 2$ smaller than the right.

\blockhead{Important words}
\begin{words}
  \item $<$ means ``less than'' (left side is smaller).
  \item Undo $-4$ by adding 4 to both sides.
\end{words}

\blockhead{Work the problem one tiny step at a time}
\begin{steps}
  \item Start with $x - 4 < 2$.
  \item Add 4 to both sides. Show under the inequality:
        \begin{center}
          $x - 4 < 2$\par\smallskip
          \underline{\hspace{0.6em}\textcolor{blue}{$+4$}\hspace{1.6em}\textcolor{blue}{$+4$}\hspace{0.6em}}
        \end{center}
  \item Left side: $x - 4 + 4$ leaves $x$.
  \item Right side leftover arithmetic: $2 + 4 = 6$.
  \item So $x < 6$.
\end{steps}

\finalans{$x < 6$}

\blockhead{Why this makes sense}
Example check: if $x = 5$, then $5 - 4 = 1$ and $1 < 2$ is true. If $x = 6$,
then $6 - 4 = 2$ and $2 < 2$ is false, so 6 is not included.

\question{H5}

\blockhead{Question}
Write an equivalent way to express $6 < y$:

\checked{$y > 6$}

\blockhead{What the question is asking}
This question asks for another comparison sentence that means the exact same
thing as $6 < y$.

\blockhead{Important words}
\begin{words}
  \item \term{Equivalent}means ``means the same thing.''
  \item To \term{flip}an inequality means swap the two sides and reverse the
        symbol ($<$ becomes $>$, or $>$ becomes $<$).
  \item If you swap the two sides of an inequality, you must also flip the
        inequality symbol.
  \item $6 < y$ means 6 is less than $y$, which is the same as saying $y$ is
        greater than 6.
\end{words}

\blockhead{Work the problem one tiny step at a time}
\begin{steps}
  \item Read $6 < y$: ``six is less than $y$.''
  \item That means $y$ is to the right of 6 on the number line.
  \item Flip both sides and flip $<$ to $>$: $y > 6$.
  \item Read $y > 6$: ``$y$ is greater than six.'' Same meaning.
\end{steps}

\finalans{$y > 6$}

\blockhead{Why this makes sense}
Do not write $y < 6$. That would reverse the meaning. Flipping sides requires
flipping the symbol.

\question{H6}

\blockhead{Question}
Graph $6 < y$ on a number line:
\blanknumline{4}{8}{1}

\checked{Open circle at 6, shade to the right (same as $y > 6$).}

\blockhead{What the question is asking}
This question asks you to show all numbers $y$ that are greater than 6 on a
straight line of numbers.

\blockhead{Important words}
\begin{words}
  \item An \term{endpoint}is the starting number on the graph where the circle
        sits.
  \item An \term{open circle}means the endpoint is \textit{not} included (used
        for $<$ or $>$).
  \item A \term{closed}(filled) circle would mean the endpoint is included (used
        for $\le$ or $\ge$).
  \item A \term{ray}on a number line is a shaded half-line that starts at a point
        and keeps going forever in one direction (shown with an arrow).
  \item Shading to the \term{right}means greater numbers. Shading to the
        \term{left}means smaller numbers.
  \item $6 < y$ means the same as $y > 6$.
\end{words}

\blockhead{Work the problem one tiny step at a time}
\begin{steps}
  \item Rewrite as $y > 6$ so $y$ is on the left (easier to graph).
  \item Find 6 on the number line.
  \item Draw an open circle at 6 because $y$ cannot equal 6.
  \item Shade the ray to the right of 6, because numbers greater than 6 are to
        the right.
  \item Put an arrow on the shaded part to show it keeps going.
        \begin{center}
        \begin{tikzpicture}[x=0.62cm,y=0.62cm]
          \draw (3.3,0) -- (8.7,0);
          \draw[-{Stealth[length=2.4mm]}] (8,0) -- (8.7,0);
          \foreach \x in {4,5,...,8}{%
            \draw (\x,0.16) -- (\x,-0.16);
            \node[below=1pt,font=\footnotesize] at (\x,-0.16) {$\x$};}
          \draw[blue,very thick,-{Stealth[length=2.6mm]}] (6,0) -- (8.6,0);
          \draw[thick,fill=white] (6,0) circle (2.6pt);
        \end{tikzpicture}
        \end{center}
\end{steps}

\finalans{Open circle at 6, shade right.}

\blockhead{Why this makes sense}
Open (not filled) circle matters. If the circle were filled, that would wrongly
include 6.

\question{H7}

\blockhead{Question}
Graph $x > 2$ on a number line:
\blanknumline{0}{5}{1}

\checked{Open circle at 2, shade to the right.}

\blockhead{What the question is asking}
This question asks you to show all numbers $x$ greater than 2 on a straight line
of numbers.

\blockhead{Important words}
\begin{words}
  \item $x > 2$ means $x$ is greater than 2 (not equal to 2).
  \item Use an \term{open circle}at 2 and shade to the \term{right}.
\end{words}

\blockhead{Work the problem one tiny step at a time}
\begin{steps}
  \item Find 2 on the number line.
  \item Draw an open circle at 2 (because $>$ does not include 2).
  \item Shade to the right of 2 for all larger numbers.
  \item Add an arrow to show the shading continues forever.
        \begin{center}
        \begin{tikzpicture}[x=0.62cm,y=0.62cm]
          \draw (-0.7,0) -- (5.7,0);
          \draw[-{Stealth[length=2.4mm]}] (5,0) -- (5.7,0);
          \foreach \x in {0,1,...,5}{%
            \draw (\x,0.16) -- (\x,-0.16);
            \node[below=1pt,font=\footnotesize] at (\x,-0.16) {$\x$};}
          \draw[blue,very thick,-{Stealth[length=2.6mm]}] (2,0) -- (5.6,0);
          \draw[thick,fill=white] (2,0) circle (2.6pt);
        \end{tikzpicture}
        \end{center}
\end{steps}

\finalans{Open circle at 2, shade right.}

\blockhead{Why this makes sense}
Numbers like 3, 4, 10, and 100 are shaded. The number 2 itself is not shaded.

\lesson{I. Geometry}

\question{I1}

\blockhead{Question}
Rectangle with length $9$ and width $4$: Area $= $

\checked{$36$ square units}

\blockhead{What the question is asking}
This question asks how much flat space is inside a rectangle with length 9 and
width 4.

\blockhead{Important words}
\begin{words}
  \item A \term{rectangle}is a four-sided shape with four right angles. Opposite
        sides are equal.
  \item \term{Length}and \term{width}are the two side lengths.
  \item \term{Area}means how much flat space is inside. For a rectangle,
        area $=$ length $\times$ width.
  \item Area is measured in \term{square units}.
\end{words}

\blockhead{Work the problem one tiny step at a time}
\begin{steps}
  \item Write the formula: area $=$ length $\times$ width.
  \item Put in the numbers: area $= 9 \times 4$.
  \item Multiply: $9 \times 4 = 36$.
  \item Include the unit idea: 36 square units.
\end{steps}

\finalans{$36$ square units}

\blockhead{Why this makes sense}
Area counts square tiles that fit inside. A 9-by-4 grid holds 36 unit squares.

\question{I2}

\blockhead{Question}
Rectangle with length $11$ and width $3$: Perimeter $= $

\checked{$28$ units}

\blockhead{What the question is asking}
This question asks for the distance all the way around a rectangle with length
11 and width 3.

\blockhead{Important words}
\begin{words}
  \item \term{Perimeter}means the distance all the way around the outside.
  \item For a rectangle, perimeter $= 2($length $+$ width$)$, or add all four
        sides.
\end{words}

\blockhead{Work the problem one tiny step at a time}
\begin{steps}
  \item Add length and width: $11 + 3 = 14$.
  \item There are two of each side, so multiply by 2: $2 \times 14 = 28$.
  \item Or add all four sides: $11 + 3 + 11 + 3 = 28$.
\end{steps}

\finalans{$28$ units}

\blockhead{Why this makes sense}
Perimeter is a length around the edge, not the space inside. Do not multiply
$11 \times 3$ here; that would be area.

\question{I3}

\blockhead{Question}
Rectangle with area $42$ sq units and width $6$: Length $= $

\checked{$7$ units}

\blockhead{What the question is asking}
This question asks for the missing length when the flat space inside and the
width are known.

\blockhead{Important words}
\begin{words}
  \item Area $=$ length $\times$ width.
  \item So length $=$ area $\div$ width.
\end{words}

\blockhead{Work the problem one tiny step at a time}
\begin{steps}
  \item Write the area equation: $42 = $ length $\times$ 6.
  \item Divide both sides by 6 to find length. Show the operation under the
        equals:
        \begin{center}
          $42 = $ length $\times$ 6\par\smallskip
          \underline{\hspace{0.6em}\textcolor{blue}{$\div 6$}\hspace{1.6em}\textcolor{blue}{$\div 6$}\hspace{0.6em}}
        \end{center}
  \item Compute the leftover division: $42 \div 6 = 7$ (because $6 \times 7 = 42$).
  \item So length $= 7$ units.
\end{steps}

\finalans{$7$ units}

\blockhead{Why this makes sense}
Check: $7 \times 6 = 42$. The area matches.

\lesson{J. Units And Quantities}

\question{J1}

\blockhead{Question}
$3$ hours $= $\ansblank\ minutes

\checked{$180$ minutes}

\blockhead{What the question is asking}
This question asks you to change 3 hours into minutes.

\blockhead{Important words}
\begin{words}
  \item A \term{conversion}changes a measurement from one unit to another.
  \item The key fact: 1 hour $=$ 60 minutes.
\end{words}

\blockhead{Work the problem one tiny step at a time}
\begin{steps}
  \item Start with 3 hours.
  \item Each hour has 60 minutes, so multiply: $3 \times 60$.
  \item $3 \times 60 = 180$.
  \item So 3 hours $=$ 180 minutes.
\end{steps}

\finalans{$180$ minutes}

\blockhead{Why this makes sense}
Three groups of 60 minutes make 180 minutes. Do not divide; hours to minutes
gets a larger number of smaller units.

\question{J2}

\blockhead{Question}
$\dfrac{3}{5}$ of $35 = $

\checked{$21$}

\blockhead{What the question is asking}
This question asks for $\dfrac{3}{5}$ of 35 (three fifths of 35).

\blockhead{Important words}
\begin{words}
  \item The word \term{of}with a fraction means multiply.
  \item $\dfrac{3}{5}$ means 3 equal parts out of 5 equal parts.
\end{words}

\blockhead{Work the problem one tiny step at a time}
\begin{steps}
  \item Write the product: $\dfrac{3}{5} \times 35$.
  \item One helpful way: first find $\dfrac{1}{5}$ of 35 by dividing
        $35 \div 5 = 7$.
  \item Then $\dfrac{3}{5}$ means 3 of those parts: $3 \times 7 = 21$.
  \item Or: $\dfrac{3 \times 35}{5} = \dfrac{105}{5} = 21$.
\end{steps}

\finalans{$21$}

\blockhead{Why this makes sense}
One fifth of 35 is 7, so three fifths is 21. That is a bit more than half of 35
(half is 17.5), and 21 fits.

\question{J3}

\blockhead{Question}
$210$ miles/$3$ hours $= $

\checked{$70$ miles per hour}

\blockhead{What the question is asking}
This question asks how many miles go with 1 hour (the amount for one hour).

\blockhead{Important words}
\begin{words}
  \item A \term{unit rate}tells how much of one quantity goes with 1 of another
        quantity.
  \item \term{Miles per hour}means miles for each 1 hour.
  \item To find a unit rate, divide the first amount by the second amount.
  \item Here: miles $\div$ hours gives miles for each 1 hour.
\end{words}

\blockhead{Work the problem one tiny step at a time}
\begin{steps}
  \item You have 210 miles and 3 hours.
  \item Divide miles by hours to get miles for 1 hour: $210 \div 3$.
  \item $3 \times 70 = 210$, so $210 \div 3 = 70$.
  \item The unit rate is 70 miles per hour.
  \item Check: 70 miles per hour $\times$ 3 hours $=$ 210 miles.
\end{steps}

\finalans{$70$ miles per hour}

\blockhead{Why this makes sense}
In 3 hours you go 210 miles, so in 1 hour you go one-third of that: 70 miles.
Multiplying back $70 \times 3 = 210$ confirms the unit rate.

\question{J4}

\blockhead{Question}
$72$ inches $= $\ansblank\ yards

\checked{$2$ yards}

\blockhead{What the question is asking}
This question asks you to change 72 inches into yards.

\blockhead{Important words}
\begin{words}
  \item A \term{unit conversion}changes a measurement from one unit name to
        another equal amount.
  \item An \term{inch}is a small length unit. A \term{yard}is a larger length
        unit.
  \item The key fact: 1 yard $=$ 36 inches.
  \item To go from a smaller unit (inches) to a larger unit (yards), divide by
        how many small units fit in one large unit.
\end{words}

\blockhead{Work the problem one tiny step at a time}
\begin{steps}
  \item Write what you know: 72 inches, and 1 yard $=$ 36 inches.
  \item Divide inches by 36: $72 \div 36$.
  \item $36 \times 2 = 72$, so $72 \div 36 = 2$.
  \item So 72 inches $=$ 2 yards.
  \item Check: 2 yards $\times$ 36 inches per yard $=$ 72 inches.
\end{steps}

\finalans{$2$ yards}

\blockhead{Why this makes sense}
Two yards means $36 + 36 = 72$ inches. Check passes: $2 \times 36 = 72$.

\question{J5}

\blockhead{Question}
$150$ miles/$5$ gallons $= $

\checked{$30$ miles per gallon}

\blockhead{What the question is asking}
This question asks how many miles go with 1 gallon.

\blockhead{Important words}
\begin{words}
  \item A \term{unit rate}tells how much of one quantity goes with 1 of another
        quantity.
  \item \term{Miles per gallon}means how many miles you travel on 1 gallon.
  \item To find a unit rate, divide the first amount by the second amount.
  \item Here: miles $\div$ gallons gives miles for each 1 gallon.
\end{words}

\blockhead{Work the problem one tiny step at a time}
\begin{steps}
  \item You have 150 miles and 5 gallons.
  \item Divide miles by gallons to get miles for 1 gallon: $150 \div 5$.
  \item $5 \times 30 = 150$, so $150 \div 5 = 30$.
  \item The unit rate is 30 miles per gallon.
  \item Check: 30 miles per gallon $\times$ 5 gallons $=$ 150 miles.
\end{steps}

\finalans{$30$ miles per gallon}

\blockhead{Why this makes sense}
Five gallons take you 150 miles, so one gallon takes you 30 miles.
